\section{Methodology}

\subsection{LinUCB}
For each arm $a$, LinUCB maintains a ridge-regression posterior

\begin{equation}
  A_a = I_d + \sum_{t:\, a_t = a} x_t x_t^\top, \qquad
  b_a = \sum_{t:\, a_t = a} r_t x_t,
  \label{eq:linucb-posterior}
\end{equation}

with point estimate $\theta_a = A_a^{-1} b_a$. The upper-confidence-bound
score for context $x$ under arm $a$ is

\begin{equation}
  p_a(x) = x^\top \theta_a + \alpha \sqrt{x^\top A_a^{-1} x},
  \label{eq:linucb-score}
\end{equation}

where the second term is an exploration bonus proportional to the model's
remaining uncertainty about $x$ under arm $a$'s posterior, and $\alpha$
controls the explore/exploit trade-off. Both $A_a$ and $b_a$ are updated in
closed form after every round; no gradient descent is involved.

\subsection{Agronomic constraint}
Let $R_a = \{(\ell_{a,j}, u_{a,j})\}_{j=1}^{7}$ be crop $a$'s training-time
feasible range (Section~3) for each raw feature $j$. At round $t$ with raw
context $\tilde{x}_t$, the feasible arm set is

\begin{equation}
  \mathcal{F}_t = \{a : \ell_{a,j} \le \tilde{x}_{t,j} \le u_{a,j} \ \forall j\}.
  \label{eq:feasible-set}
\end{equation}

\emph{Constrained} LinUCB selects $\arg\max_{a \in \mathcal{F}_t} p_a(x_t)$;
plain LinUCB selects over every arm. If $\mathcal{F}_t = \emptyset$, every
arm is treated as feasible rather than leaving the policy with no action.
